% roomscan — SHORT version, Thai — the Thai companion of main_short.tex (English); chapters are added as
% they are translated (so far: บทคัดย่อ, บทนำ, ทฤษฎีและงานวิจัยที่เกี่ยวข้อง).
% Same macros (shortmacros.tex), numbers (numbers.tex) and references (refs.bib) as the English version,
% so the two stay in sync when numbers change.
% Build:  cd paper/latex && tectonic main_short_th.tex   (needs XeTeX: Thai font + Thai line breaking)
\documentclass[conference]{IEEEtran}
\IEEEoverridecommandlockouts

\usepackage{fontspec}
% Thai main font: the first one installed wins (Windows / macOS / Linux). Latin text uses its Latin glyphs.
\IfFontExistsTF{TH Sarabun New}{\setmainfont{TH Sarabun New}[Script=Thai,Scale=1.35]}{%
\IfFontExistsTF{Sarabun}{\setmainfont{Sarabun}[Script=Thai]}{%
\IfFontExistsTF{Leelawadee UI}{\setmainfont{Leelawadee UI}[Script=Thai]}{%
\IfFontExistsTF{Thonburi}{\setmainfont{Thonburi}[Script=Thai]}{%
\IfFontExistsTF{Tahoma}{\setmainfont{Tahoma}[Script=Thai]}{%
\setmainfont{Noto Serif Thai}[Script=Thai]}}}}}
\setmonofont{texgyrecursor}[Extension=.otf,UprightFont=*-regular,BoldFont=*-bold,
  ItalicFont=*-italic,BoldItalicFont=*-bolditalic]
% Thai has no spaces between words: break by dictionary, with only a little stretch at those breaks so
% justification widens real spaces rather than opening gaps inside words.
\XeTeXlinebreaklocale "th"
\XeTeXlinebreakskip=0pt plus 0.1pt

\usepackage{amsmath,amssymb}
\usepackage{url}
\urlstyle{tt}
\usepackage{siunitx}
\usepackage[hidelinks]{hyperref}
\usepackage{cite}

% Macros shared by main_short.tex (English) and main_short_th.tex (Thai). Load after siunitx.
\sisetup{detect-all,range-phrase=--,range-units=single,separate-uncertainty=true,multi-part-units=single}
\newcommand{\code}[1]{\texttt{#1}}
\newcommand{\fscore}{F@\SI{5}{cm}}
% The three fine-tunes, named teacher + number of training rooms (runs mono_ft_faro / mono_ft_lidar /
% mono_ft_lidar_all; ADR-014 calls them R1 / R2 / R3). \ftmain is the paper's main model.
\newcommand{\ftfaro}{\mbox{FT-Faro-10}}
\newcommand{\ftlidar}{\mbox{FT-LiDAR-10}}
\newcommand{\ftmain}{\mbox{FT-LiDAR-24}}
\hyphenation{LiDAR}   % never "Li-DAR"

\renewcommand{\abstractname}{บทคัดย่อ}
\renewcommand{\IEEEkeywordsname}{คำสำคัญ}
\renewcommand{\refname}{เอกสารอ้างอิง}
% IEEEtran sets the abstract and keywords in bold; a bold Thai paragraph is hard to read, so use regular.
\newenvironment{thaiabstract}{\begin{abstract}\mdseries}{\end{abstract}}

\begin{document}

\title{สอนหน่วยเมตรให้โมเดลความลึกจากภาพเดี่ยวด้วย LiDAR ของโทรศัพท์เอง:\\
เปรียบเทียบ Depth Anything~V2 แบบ fine-tune กับแบบ pretrained สำหรับการสร้างแบบจำลองห้อง
โดยวัดเทียบกับ laser scanner}

% TODO: fill in author name(s) and affiliation before submission
\author{\IEEEauthorblockN{Author Name}
\IEEEauthorblockA{Affiliation\\
email@example.com}}

\maketitle

% generated by scripts/paper_numbers.py — do not edit

\newcommand{\rsGTCM}{2.1}
\newcommand{\rsGTSD}{0.5}
\newcommand{\rsGTF}{0.97}
\newcommand{\rsGTABSREL}{0.000}
\newcommand{\rsGTACC}{1.3}
\newcommand{\rsGTCOMP}{2.8}
\newcommand{\rsGTFtwo}{0.86}
\newcommand{\rsGTRtwo}{0.84}
\newcommand{\rsGTFfive}{0.97}
\newcommand{\rsGTRfive}{0.96}
\newcommand{\rsGTFten}{0.99}
\newcommand{\rsGTRten}{0.98}
\newcommand{\rsLIDARCM}{3.0}
\newcommand{\rsLIDARSD}{0.6}
\newcommand{\rsLIDARF}{0.93}
\newcommand{\rsLIDARABSREL}{0.018}
\newcommand{\rsLIDARACC}{2.2}
\newcommand{\rsLIDARCOMP}{3.9}
\newcommand{\rsLIDARFtwo}{0.54}
\newcommand{\rsLIDARRtwo}{0.52}
\newcommand{\rsLIDARFfive}{0.93}
\newcommand{\rsLIDARRfive}{0.91}
\newcommand{\rsLIDARFten}{0.98}
\newcommand{\rsLIDARRten}{0.97}
\newcommand{\rsORACLECM}{5.3}
\newcommand{\rsORACLESD}{1.3}
\newcommand{\rsORACLEF}{0.79}
\newcommand{\rsORACLEABSREL}{0.032}
\newcommand{\rsORACLEACC}{5.9}
\newcommand{\rsORACLECOMP}{4.7}
\newcommand{\rsORACLEFtwo}{0.46}
\newcommand{\rsORACLERtwo}{0.47}
\newcommand{\rsORACLEFfive}{0.79}
\newcommand{\rsORACLERfive}{0.82}
\newcommand{\rsORACLEFten}{0.90}
\newcommand{\rsORACLERten}{0.93}
\newcommand{\rsSPARSECM}{6.0}
\newcommand{\rsSPARSESD}{1.7}
\newcommand{\rsSPARSEF}{0.74}
\newcommand{\rsSPARSEABSREL}{0.036}
\newcommand{\rsSPARSEACC}{6.4}
\newcommand{\rsSPARSECOMP}{5.5}
\newcommand{\rsSPARSEFtwo}{0.38}
\newcommand{\rsSPARSERtwo}{0.39}
\newcommand{\rsSPARSEFfive}{0.74}
\newcommand{\rsSPARSERfive}{0.77}
\newcommand{\rsSPARSEFten}{0.88}
\newcommand{\rsSPARSERten}{0.91}
\newcommand{\rsSCENECM}{22.9}
\newcommand{\rsSCENESD}{2.1}
\newcommand{\rsSCENEF}{0.23}
\newcommand{\rsSCENEABSREL}{0.237}
\newcommand{\rsSCENEACC}{24.2}
\newcommand{\rsSCENECOMP}{21.6}
\newcommand{\rsSCENEFtwo}{0.08}
\newcommand{\rsSCENERtwo}{0.09}
\newcommand{\rsSCENEFfive}{0.23}
\newcommand{\rsSCENERfive}{0.27}
\newcommand{\rsSCENEFten}{0.41}
\newcommand{\rsSCENERten}{0.47}
\newcommand{\rsMETRICCM}{54.3}
\newcommand{\rsMETRICSD}{10.6}
\newcommand{\rsMETRICF}{0.03}
\newcommand{\rsMETRICABSREL}{0.382}
\newcommand{\rsMETRICACC}{70.3}
\newcommand{\rsMETRICCOMP}{38.3}
\newcommand{\rsMETRICFtwo}{0.01}
\newcommand{\rsMETRICRtwo}{0.01}
\newcommand{\rsMETRICFfive}{0.03}
\newcommand{\rsMETRICRfive}{0.04}
\newcommand{\rsMETRICFten}{0.06}
\newcommand{\rsMETRICRten}{0.10}
\newcommand{\rsGTa}{1.8}
\newcommand{\rsGTb}{1.2}
\newcommand{\rsGTc}{1.9}
\newcommand{\rsGTd}{2.8}
\newcommand{\rsGTe}{2.4}
\newcommand{\rsGTf}{2.3}
\newcommand{\rsLIDARa}{2.8}
\newcommand{\rsLIDARb}{2.4}
\newcommand{\rsLIDARc}{2.9}
\newcommand{\rsLIDARd}{4.1}
\newcommand{\rsLIDARe}{3.1}
\newcommand{\rsLIDARf}{3.0}
\newcommand{\rsORACLEa}{3.8}
\newcommand{\rsORACLEb}{6.2}
\newcommand{\rsORACLEc}{4.8}
\newcommand{\rsORACLEd}{7.5}
\newcommand{\rsORACLEe}{5.2}
\newcommand{\rsORACLEf}{4.5}
\newcommand{\rsSPARSEa}{4.5}
\newcommand{\rsSPARSEb}{6.1}
\newcommand{\rsSPARSEc}{5.5}
\newcommand{\rsSPARSEd}{9.1}
\newcommand{\rsSPARSEe}{5.9}
\newcommand{\rsSPARSEf}{4.6}
\newcommand{\rsSCENEa}{24.1}
\newcommand{\rsSCENEb}{26.0}
\newcommand{\rsSCENEc}{19.9}
\newcommand{\rsSCENEd}{23.2}
\newcommand{\rsSCENEe}{21.4}
\newcommand{\rsSCENEf}{23.1}
\newcommand{\rsMETRICa}{48.5}
\newcommand{\rsMETRICb}{61.6}
\newcommand{\rsMETRICc}{41.5}
\newcommand{\rsMETRICd}{48.9}
\newcommand{\rsMETRICe}{54.1}
\newcommand{\rsMETRICf}{71.0}
\newcommand{\rsFTFAROa}{10.8}
\newcommand{\rsFTFAROb}{13.8}
\newcommand{\rsFTFAROc}{17.6}
\newcommand{\rsFTFAROd}{21.1}
\newcommand{\rsFTFAROe}{12.4}
\newcommand{\rsFTFAROf}{11.9}
\newcommand{\rsFTLIDARa}{13.9}
\newcommand{\rsFTLIDARb}{19.0}
\newcommand{\rsFTLIDARc}{12.8}
\newcommand{\rsFTLIDARd}{12.5}
\newcommand{\rsFTLIDARe}{16.6}
\newcommand{\rsFTLIDARf}{24.6}
\newcommand{\rsFTLIDARALLa}{12.8}
\newcommand{\rsFTLIDARALLb}{16.7}
\newcommand{\rsFTLIDARALLc}{17.3}
\newcommand{\rsFTLIDARALLd}{12.2}
\newcommand{\rsFTLIDARALLe}{12.4}
\newcommand{\rsFTLIDARALLf}{22.5}
\newcommand{\rsSPARSEGAP}{0.6}
\newcommand{\rsSPARSEGAPLIDAR}{2.9}
\newcommand{\rsSPARSEVSSCENE}{3.8}
\newcommand{\rsLIDARGAPGTCM}{0.98}
\newcommand{\rsLIDARGAPGTSD}{0.25}
\newcommand{\rsORACLEGAPLIDARCM}{2.29}
\newcommand{\rsORACLEGAPLIDARSD}{1.11}
\newcommand{\rsCONFONECM}{3.07}
\newcommand{\rsCONFONEF}{0.93}
\newcommand{\rsCONFONEACC}{2.17}
\newcommand{\rsCONFONECOMP}{3.98}
\newcommand{\rsCONFTWOCM}{3.66}
\newcommand{\rsCONFTWOF}{0.93}
\newcommand{\rsCONFTWOACC}{2.08}
\newcommand{\rsCONFTWOCOMP}{5.23}
\newcommand{\rsWZOTCM}{3.07}
\newcommand{\rsWZOTF}{0.93}
\newcommand{\rsWZOTACC}{2.17}
\newcommand{\rsWZOTCOMP}{3.98}
\newcommand{\rsWOTFCM}{3.04}
\newcommand{\rsWOTFF}{0.93}
\newcommand{\rsWOTFACC}{2.23}
\newcommand{\rsWOTFCOMP}{3.85}
\newcommand{\rsFTFAROCM}{14.6}
\newcommand{\rsFTFAROSD}{4.0}
\newcommand{\rsFTFAROF}{0.31}
\newcommand{\rsFTFAROABSREL}{0.136}
\newcommand{\rsFTFARODELTAONE}{0.780}
\newcommand{\rsFTFARORtwo}{0.12}
\newcommand{\rsFTFARORfive}{0.34}
\newcommand{\rsFTFARORten}{0.57}
\newcommand{\rsFTLIDARCM}{16.6}
\newcommand{\rsFTLIDARSD}{4.7}
\newcommand{\rsFTLIDARF}{0.23}
\newcommand{\rsFTLIDARABSREL}{0.136}
\newcommand{\rsFTLIDARDELTAONE}{0.848}
\newcommand{\rsFTLIDARRtwo}{0.09}
\newcommand{\rsFTLIDARRfive}{0.26}
\newcommand{\rsFTLIDARRten}{0.49}
\newcommand{\rsFTLIDARALLCM}{15.7}
\newcommand{\rsFTLIDARALLSD}{4.0}
\newcommand{\rsFTLIDARALLF}{0.24}
\newcommand{\rsFTLIDARALLABSREL}{0.134}
\newcommand{\rsFTLIDARALLDELTAONE}{0.812}
\newcommand{\rsFTLIDARALLRtwo}{0.09}
\newcommand{\rsFTLIDARALLRfive}{0.28}
\newcommand{\rsFTLIDARALLRten}{0.51}
\newcommand{\rsFTTEACHERRATIO}{1.13}
\newcommand{\rsFTGAIN}{3.3}
\newcommand{\rsFTGAPORACLE}{11.2}
\newcommand{\rsDALARGECM}{6.45}
\newcommand{\rsDALARGESD}{2.42}
\newcommand{\rsDALARGEF}{0.74}
\newcommand{\rsDALARGEABSREL}{0.040}
\newcommand{\rsDALARGEDELTAONE}{0.979}
\newcommand{\rsDALARGESCAN}{355}
\newcommand{\rsDASMALLCM}{9.22}
\newcommand{\rsDASMALLSD}{3.48}
\newcommand{\rsDASMALLF}{0.64}
\newcommand{\rsDASMALLABSREL}{0.054}
\newcommand{\rsDASMALLDELTAONE}{0.964}
\newcommand{\rsDASMALLSCAN}{56}
\newcommand{\rsMIDASCM}{14.95}
\newcommand{\rsMIDASSD}{1.87}
\newcommand{\rsMIDASF}{0.35}
\newcommand{\rsMIDASABSREL}{0.110}
\newcommand{\rsMIDASDELTAONE}{0.873}
\newcommand{\rsMIDASSCAN}{18}
\newcommand{\rsSTRGTONECM}{2.1}
\newcommand{\rsSTRGTONEACC}{1.3}
\newcommand{\rsSTRGTONECOMP}{2.8}
\newcommand{\rsSTRGTONEF}{0.97}
\newcommand{\rsSTRGTONEABSREL}{0.000}
\newcommand{\rsSTRGTONEN}{391}
\newcommand{\rsSTRGTONESCAN}{2}
\newcommand{\rsSTRGTFIVECM}{4.4}
\newcommand{\rsSTRGTFIVEACC}{1.2}
\newcommand{\rsSTRGTFIVECOMP}{7.5}
\newcommand{\rsSTRGTFIVEF}{0.91}
\newcommand{\rsSTRGTFIVEABSREL}{0.000}
\newcommand{\rsSTRGTFIVEN}{78}
\newcommand{\rsSTRGTFIVESCAN}{1}
\newcommand{\rsSTRGTTENCM}{8.9}
\newcommand{\rsSTRGTTENACC}{1.2}
\newcommand{\rsSTRGTTENCOMP}{16.6}
\newcommand{\rsSTRGTTENF}{0.83}
\newcommand{\rsSTRGTTENABSREL}{0.000}
\newcommand{\rsSTRGTTENN}{40}
\newcommand{\rsSTRGTTENSCAN}{1}
\newcommand{\rsSTRGTTWENTYCM}{16.8}
\newcommand{\rsSTRGTTWENTYACC}{1.2}
\newcommand{\rsSTRGTTWENTYCOMP}{32.3}
\newcommand{\rsSTRGTTWENTYF}{0.68}
\newcommand{\rsSTRGTTWENTYABSREL}{0.000}
\newcommand{\rsSTRGTTWENTYN}{20}
\newcommand{\rsSTRGTTWENTYSCAN}{1}
\newcommand{\rsSTRORACLEONECM}{6.5}
\newcommand{\rsSTRORACLEONEACC}{7.1}
\newcommand{\rsSTRORACLEONECOMP}{5.8}
\newcommand{\rsSTRORACLEONEF}{0.74}
\newcommand{\rsSTRORACLEONEABSREL}{0.040}
\newcommand{\rsSTRORACLEONEN}{391}
\newcommand{\rsSTRORACLEONESCAN}{358}
\newcommand{\rsSTRORACLEFIVECM}{8.2}
\newcommand{\rsSTRORACLEFIVEACC}{5.6}
\newcommand{\rsSTRORACLEFIVECOMP}{10.9}
\newcommand{\rsSTRORACLEFIVEF}{0.72}
\newcommand{\rsSTRORACLEFIVEABSREL}{0.040}
\newcommand{\rsSTRORACLEFIVEN}{78}
\newcommand{\rsSTRORACLEFIVESCAN}{72}
\newcommand{\rsSTRORACLETENCM}{12.7}
\newcommand{\rsSTRORACLETENACC}{5.0}
\newcommand{\rsSTRORACLETENCOMP}{20.3}
\newcommand{\rsSTRORACLETENF}{0.65}
\newcommand{\rsSTRORACLETENABSREL}{0.041}
\newcommand{\rsSTRORACLETENN}{40}
\newcommand{\rsSTRORACLETENSCAN}{36}
\newcommand{\rsSTRORACLETWENTYCM}{20.9}
\newcommand{\rsSTRORACLETWENTYACC}{5.1}
\newcommand{\rsSTRORACLETWENTYCOMP}{36.7}
\newcommand{\rsSTRORACLETWENTYF}{0.52}
\newcommand{\rsSTRORACLETWENTYABSREL}{0.040}
\newcommand{\rsSTRORACLETWENTYN}{20}
\newcommand{\rsSTRORACLETWENTYSCAN}{19}
\newcommand{\rsSTRSCENEONECM}{22.7}
\newcommand{\rsSTRSCENEONEACC}{23.8}
\newcommand{\rsSTRSCENEONECOMP}{21.7}
\newcommand{\rsSTRSCENEONEF}{0.22}
\newcommand{\rsSTRSCENEONEABSREL}{0.240}
\newcommand{\rsSTRSCENEONEN}{391}
\newcommand{\rsSTRSCENEONESCAN}{346}
\newcommand{\rsSTRSCENEFIVECM}{22.7}
\newcommand{\rsSTRSCENEFIVEACC}{20.2}
\newcommand{\rsSTRSCENEFIVECOMP}{25.2}
\newcommand{\rsSTRSCENEFIVEF}{0.23}
\newcommand{\rsSTRSCENEFIVEABSREL}{0.236}
\newcommand{\rsSTRSCENEFIVEN}{78}
\newcommand{\rsSTRSCENEFIVESCAN}{71}
\newcommand{\rsSTRSCENETENCM}{25.0}
\newcommand{\rsSTRSCENETENACC}{18.9}
\newcommand{\rsSTRSCENETENCOMP}{31.0}
\newcommand{\rsSTRSCENETENF}{0.20}
\newcommand{\rsSTRSCENETENABSREL}{0.247}
\newcommand{\rsSTRSCENETENN}{40}
\newcommand{\rsSTRSCENETENSCAN}{36}
\newcommand{\rsSTRSCENETWENTYCM}{34.5}
\newcommand{\rsSTRSCENETWENTYACC}{21.4}
\newcommand{\rsSTRSCENETWENTYCOMP}{47.6}
\newcommand{\rsSTRSCENETWENTYF}{0.14}
\newcommand{\rsSTRSCENETWENTYABSREL}{0.268}
\newcommand{\rsSTRSCENETWENTYN}{20}
\newcommand{\rsSTRSCENETWENTYSCAN}{19}

% บทคัดย่อภาษาไทย ใช้ใน main_short_th.tex (ฉบับอังกฤษอยู่ใน 00_abstract.tex); ตัวเลขทั้งหมดมาจาก numbers.tex.
% environment thaiabstract นิยามใน main_short_th.tex; ต้อง build ด้วย XeTeX/tectonic.
\begin{thaiabstract}
โทรศัพท์ส่วนใหญ่ไม่มี LiDAR จึงต้องสแกนห้องด้วยการประมาณความลึกจากภาพเดี่ยว
แต่โมเดลสำเร็จรูปที่ให้ความลึกเป็นเมตร เช่น Depth Anything V2 (DA-v2) Metric-Indoor~\cite{da2}
ซึ่งฝึกด้วยภาพสังเคราะห์ ทำให้ห้องที่สร้างได้คลาดเคลื่อน \rsMETRICCM{}~ซม.
งานวิจัยนี้ศึกษาว่าการปรับจูน (fine-tune) โมเดลดังกล่าวด้วยความลึกจาก LiDAR ของโทรศัพท์เอง
โดยไม่ใช้เครื่องสแกนเลเซอร์ ทำให้แม่นยำกว่าโมเดลเดิมหรือไม่ เมื่อวัดเทียบกับเครื่องสแกนเลเซอร์
Faro Focus S70~\cite{faro} ซึ่งเป็นค่าอ้างอิงของชุดข้อมูล ARKitScenes~\cite{arkitscenes}
โดยปรับจูนเฉพาะส่วน DPT head~\cite{dpt} (ตัวถอดรหัสความลึก โดยตรึงส่วนเข้ารหัสภาพไว้)
ด้วยความลึกจาก LiDAR ของ iPad บน 24 ห้อง ได้โมเดล \ftmain{}
แล้วเปรียบเทียบกับโมเดลเดิมในกระบวนการเดียวกัน คือรวมความลึกรายเฟรมเป็นปริมาตร
TSDF (truncated signed distance function)~\cite{curless96} แล้วสร้างเป็นพื้นผิวสามมิติ
โดยให้แหล่งความลึกเป็นตัวแปรเดียว บนห้องที่ไม่เคยใช้ฝึก ได้แก่ 6 ห้องที่มีแบบจำลองอ้างอิงจากเครื่องสแกนเลเซอร์
และอีก \rsVN{} ห้องที่ประเมินรายเฟรม
ผลคือ \ftmain{} ลดค่า Chamfer distance~\cite{chamfer}
(ระยะเฉลี่ยถึงจุดที่ใกล้ที่สุดระหว่างพื้นผิวที่สร้างได้กับพื้นผิวอ้างอิง วัดทั้งสองทิศ) จาก \rsMETRICCM{}~ซม. เหลือ \rsFTLIDARALLCM{}~ซม.
(ดีขึ้น \rsMAINGAIN{} เท่า และดีขึ้นทุกห้อง) และลดค่า AbsRel~\cite{eigen14}
(ความคลาดเคลื่อนสัมพัทธ์เฉลี่ยของความลึก) จาก \rsVMETRICABSREL{} เหลือ \rsVFTALLABSREL{}
(ดีขึ้น \rsVWINS{} จาก \rsVN{} ห้อง) โดยไม่ต้องปรับเทียบหรือใช้เซนเซอร์ขณะใช้งาน
และการใช้เครื่องสแกนเลเซอร์เป็นครูไม่ได้ดีกว่าอย่างวัดได้
ความคลาดเคลื่อนที่เหลือมาจากสเกลที่แกว่งตามมุมมองของแต่ละเฟรม ซึ่งการปรับจูนขจัดไม่ได้
จึงยังด้อยกว่า LiDAR ของ iPad (\rsLIDARCM{}~ซม.) และยังไม่ถึงระดับความแม่นยำสำหรับดูภาพรวม
(พื้นผิวของห้องอย่างน้อย 90\,\% อยู่ภายใน 10~ซม.)
\end{thaiabstract}


\begin{IEEEkeywords}\mdseries
การประมาณความลึกจากภาพเดี่ยว (monocular depth estimation), การเรียนรู้แบบถ่ายโอน (transfer learning),
ความลึกแบบเมตริก (metric depth), การใช้ LiDAR เป็นข้อมูลสอน (LiDAR supervision), การรวมความลึกแบบ TSDF
(TSDF fusion), การสร้างแบบจำลองสามมิติของห้อง (3D room reconstruction), ARKitScenes
\end{IEEEkeywords}

% บทนำภาษาไทย ใช้ใน main_short_th.tex (ฉบับอังกฤษอยู่ใน 01_introduction.tex); แก้ให้ตรงกันทั้งสองฉบับ
% ตัวเลขทั้งหมดมาจาก numbers.tex. ไม่ใช้ \emph เพราะฟอนต์ไทยไม่มีตัวเอียง.
\section{บทนำ}\label{sec:intro}

การสแกนห้องเป็นแบบจำลองสามมิติเป็นความสามารถที่มีอยู่แล้วใน iPhone~Pro และ iPad~Pro
(เช่น Apple RoomPlan~\cite{roomplan}) โดยอาศัยเซนเซอร์ LiDAR วัดความลึกเป็นหน่วยเมตร
แต่อุปกรณ์ที่มี LiDAR เป็นส่วนน้อย iPhone รุ่นที่ไม่ใช่ Pro และโทรศัพท์ Android แทบทุกรุ่นไม่มีเซนเซอร์นี้
ทางเลือกของอุปกรณ์เหล่านี้คือการประมาณความลึกจากภาพเดี่ยว (monocular depth estimation)
คือการทำนายความลึกจากภาพสีเพียงภาพเดียว ซึ่งพัฒนาอย่างรวดเร็วตั้งแต่ MiDaS~\cite{midas}
จนถึง Depth Anything~\cite{da1,da2} ทีมพัฒนาผลิตภัณฑ์จึงต้องรู้ว่า
\textbf{หากใช้โมเดลจากภาพเดี่ยวแทน LiDAR ห้องที่ได้จะคลาดเคลื่อนกี่เซนติเมตร และยังใช้งานแบบใดได้บ้าง}
เกณฑ์วัดมาตรฐานแบบรายภาพ~\cite{eigen14} ซึ่งเฉลี่ยความคลาดเคลื่อนของความลึกบนพิกเซลของภาพแต่ละภาพ
ตอบคำถามนี้ไม่ได้ และระบบสร้างแบบจำลองแบบครบวงจร (end-to-end)~\cite{atlas,neuralrecon,simplerecon}
ก็ตอบไม่ได้ เพราะออกแบบกระบวนการทั้งหมดรอบโมเดลของตนเอง

โมเดลรุ่นใหม่บางตัวให้ความลึกเป็นเมตรได้โดยตรง~\cite{da2,zoedepth,depthpro}
แต่ฝึกด้วยภาพที่ไม่ได้มาจากกล้องโทรศัพท์ เช่น \mbox{DA-v2} Metric-Indoor ฝึกด้วยภาพห้องสังเคราะห์จากชุดข้อมูล
Hypersim~\cite{hypersim} และในการทดลองของเรา ห้องที่สร้างด้วย \mbox{DA-v2} Metric-Indoor
คลาดเคลื่อนเฉลี่ย \rsMETRICCM{}~ซม. วิธีแก้โดยตรงคือการปรับจูน (fine-tune)
โมเดลด้วยความลึกจากอุปกรณ์เป้าหมาย คำถามคือข้อมูลสอน (label) นี้จะมาจากไหน
ระหว่างเครื่องสแกนเลเซอร์ระดับงานสำรวจ~\cite{faro} ซึ่งผู้พัฒนาแอปเพียงไม่กี่รายมีใช้
กับ LiDAR ที่มีอยู่แล้วในอุปกรณ์รุ่น Pro ซึ่งแอปเก็บความลึกได้ระหว่างการใช้งานปกติ
หากวิธีหลังได้ผล อุปกรณ์ที่มี LiDAR ก็จะเป็นครูสอนอุปกรณ์ที่ไม่มี LiDAR ได้

งานวิจัยนี้จึงตั้งคำถามว่า \mbox{DA-v2} Metric-Indoor ที่ปรับจูนด้วยความลึกจาก LiDAR ของ iPad
โดยไม่ใช้เครื่องสแกนเลเซอร์ในการฝึก แม่นยำกว่าโมเดลเดิมที่ไม่ได้ปรับจูนหรือไม่
เมื่อวัดเทียบกับเครื่องสแกนเลเซอร์ Faro ซึ่งเป็นค่าอ้างอิงของชุดข้อมูล ARKitScenes~\cite{arkitscenes}
บนห้องที่ไม่เคยใช้ฝึก โดยมีวัตถุประสงค์ 4 ข้อ ได้แก่
(1)~ปรับจูนโมเดลด้วยความลึกจาก LiDAR ของ iPad จำนวน 24 ห้อง ได้โมเดล \ftmain{}
แล้วเปรียบเทียบกับโมเดลเดิมบนห้องที่ไม่เคยใช้ฝึก โดยถือว่าดีกว่าเมื่อค่าเฉลี่ยความคลาดเคลื่อนต่ำกว่า
และผลต่างรายห้องมีนัยสำคัญทางสถิติที่ระดับ 5\,\% ตามการทดสอบ Wilcoxon signed-rank แบบ exact~\cite{wilcoxon}
(2)~วัดว่าแหล่งของข้อมูลสอนมีผลหรือไม่ โดยเปรียบเทียบข้อมูลสอนจาก Faro กับจาก LiDAR
บน 10 ห้องเดียวกัน (\ftfaro{} กับ \ftlidar{}) เปรียบเทียบข้อมูลสอนจาก LiDAR 10 ห้องกับ 24 ห้อง
และเปรียบเทียบกับโมเดลที่ใช้ความยาวโฟกัสของกล้องแทนข้อมูลสอนจากโดเมนเป้าหมาย (Depth Pro~\cite{depthpro})
(3)~ระบุตำแหน่งของ \ftmain{} เทียบกับแหล่งความลึกที่ใช้ได้ขณะใช้งาน ได้แก่ LiDAR ของ iPad
และโมเดลความลึกสัมพัทธ์ที่ปรับสเกลรายเฟรมด้วยจุดเบาบาง (sparse points) หรือปรับสเกลครั้งเดียวต่อห้อง
และ (4)~หาที่มาของความคลาดเคลื่อนที่การปรับจูนแก้ไม่ได้ โดยหาสเกลความลึกที่เหมาะกับแต่ละเฟรม
แล้วทดสอบว่าสเกลนั้นเปลี่ยนตามเวลา ซึ่งบ่งชี้การสะสมความคลาดเคลื่อน (drift) หรือเปลี่ยนตามสิ่งที่อยู่ในภาพ

ทุกการเปรียบเทียบวัดสองแบบ แบบแรกวัดเป็นแบบจำลองห้อง (mesh) ด้วย Chamfer distance~\cite{chamfer}
คือระยะเฉลี่ยจากแต่ละพื้นผิวถึงจุดที่ใกล้ที่สุดของอีกพื้นผิว และด้วย F-score ที่ 5~ซม.
(\fscore{})~\cite{tanks} ซึ่งรวมสัดส่วนของพื้นผิวที่สร้างได้ที่อยู่ภายใน 5~ซม. จากพื้นผิวอ้างอิง
กับสัดส่วนของพื้นผิวอ้างอิงที่ถูกสร้างขึ้นภายใน 5~ซม. แบบที่สองวัดรายเฟรมด้วย AbsRel~\cite{eigen14}
คือความคลาดเคลื่อนสัมพัทธ์เฉลี่ยของความลึก และด้วยสัดส่วนของพิกเซลที่ความลึกที่ทำนายต่างจากค่าจริงไม่เกิน 1.25 เท่า
(ไม่ว่าจะมากหรือน้อยกว่าค่าจริง) ซึ่งเรียกว่า \wone{}~\cite{eigen14}
สำหรับวัตถุประสงค์ข้อ (3) สัดส่วนของพื้นผิวอ้างอิงที่ถูกสร้างขึ้นภายใน 2, 5 และ 10~ซม.
ใช้แทนความแม่นยำที่งานวัดเพื่อปรับปรุงห้อง งานวางเฟอร์นิเจอร์ และงานดูภาพรวมของห้องต้องการ ตามลำดับ
และถือว่ารองรับงานนั้นได้เมื่อสัดส่วนนี้ถึง 90\,\%

ขอบเขตของงานวิจัยมีดังนี้ ข้อมูลเป็นห้องในบ้านจากชุดข้อมูล ARKitScenes ที่ถ่ายด้วย iPad~Pro รุ่นเดียว
(ภาพสีขนาด $640\times480$ พิกเซล และความลึกจาก LiDAR ขนาด $256\times192$ พิกเซล)
ค่าอ้างอิงคือความลึกจาก Faro ที่มากับชุดข้อมูล และไม่ใช้ LiDAR ของ iPad เป็นค่าอ้างอิง
การปรับจูนทำเฉพาะส่วน DPT head~\cite{dpt} ซึ่งเป็นตัวถอดรหัสความลึกของ \mbox{DA-v2} Metric-Indoor ขนาด Large
ด้วยสูตรการฝึกเดียว โดยตรึงส่วนเข้ารหัสภาพไว้ ประเมินผลบนห้องทดสอบ 6 ห้องทั้งแบบแบบจำลองห้องและแบบรายเฟรม
และอีก \rsVN{} ห้องแบบรายเฟรมเท่านั้น ทุกห้องไม่ซ้ำกับห้องที่ใช้ฝึกหรือใช้เลือก checkpoint
ทุกแหล่งความลึกผ่านกระบวนการแบบออฟไลน์เดียวกัน คือความลึกรายเฟรม การปรับสเกล
การรวมเป็นปริมาตร TSDF (truncated signed distance function)~\cite{curless96} และการสร้างพื้นผิว
แหล่งความลึกจึงเป็นตัวแปรเดียว นอกจากนี้ ทุกแหล่งและค่าอ้างอิงใช้ตำแหน่งกล้องชุดเดียวกัน
ซึ่งบันทึกโดย ARKit ความคลาดเคลื่อนของตำแหน่งกล้องจึงไม่อยู่ในผลการวัด สิ่งที่อยู่นอกขอบเขต ได้แก่
ระบบสร้างแบบจำลองที่เรียนรู้ทั้งกระบวนการ~\cite{atlas,neuralrecon,simplerecon}
กล้องโทรศัพท์รุ่นอื่นรวมถึง Android การประมวลผลแบบเรียลไทม์บนอุปกรณ์ และการประมาณตำแหน่งกล้อง

ประโยชน์ที่คาดว่าจะได้รับสอดคล้องกับวัตถุประสงค์ทีละข้อ ข้อแรก ผู้พัฒนาแอปจะรู้ว่าข้อมูลที่เก็บได้อยู่แล้ว
จากผู้ใช้อุปกรณ์รุ่น Pro ช่วยให้การสแกนห้องบนอุปกรณ์ที่ไม่มี LiDAR ดีขึ้นหรือไม่
โดยไม่ต้องปรับเทียบหรือใช้เซนเซอร์เพิ่มขณะใช้งาน ข้อสอง จะรู้ว่าการสร้างข้อมูลสอนจำเป็นต้องใช้เครื่องสแกนเลเซอร์หรือไม่
ข้อสาม จะรู้ว่าอุปกรณ์ที่ไม่มี LiDAR รองรับงานแบบใดได้บ้าง
พร้อมได้กระบวนการที่ประเมินแหล่งความลึกใหม่ได้ด้วยการแก้ค่าตั้งค่าเพียงรายการเดียว
และข้อสี่ จะรู้ว่าควรลงทุนพัฒนาส่วนใดต่อ ระหว่างสเกลรายเฟรม ข้อมูลฝึก หรือขนาดโมเดล

% ทฤษฎีและงานวิจัยที่เกี่ยวข้อง ภาษาไทย ใช้ใน main_short_th.tex (ฉบับอังกฤษอยู่ใน 02_related_work.tex);
% แก้ให้ตรงกันทั้งสองฉบับ. ไม่ใช้ \emph เพราะฟอนต์ไทยไม่มีตัวเอียง.
\section{ทฤษฎีและงานวิจัยที่เกี่ยวข้อง}\label{sec:related}

บทนี้อธิบายทฤษฎีที่งานวิจัยนี้อาศัย ได้แก่ เหตุใดภาพเดียวจึงบอกความลึกเป็นเมตรไม่ได้
โมเดลที่เรานำมาปรับจูนมีโครงสร้างและฝึกมาอย่างไร และความลึกรายภาพกลายเป็นแบบจำลองห้องได้อย่างไร
พร้อมทั้งวางตำแหน่งของงานนี้เทียบกับงานก่อนหน้าในแต่ละประเด็น

\subsection{การประมาณความลึกจากภาพเดี่ยว}
ตามแบบจำลองกล้องรูเข็ม (pinhole camera)~\cite{hartley} ขนาดของวัตถุในภาพแปรผันตามขนาดจริงของวัตถุ
หารด้วยระยะห่าง แล้วคูณด้วยความยาวโฟกัสของเลนส์ ห้องที่ใหญ่เป็นสองเท่าเมื่อมองจากระยะไกลเป็นสองเท่า
จึงให้ภาพเหมือนกันทุกประการ เรขาคณิตเพียงอย่างเดียวจึงหาความลึกเป็นเมตรจากภาพเดียวไม่ได้
โครงข่ายประสาทเทียมทำได้เพียงอนุมานจากสิ่งที่เรียนรู้มา เช่น ความสูงปกติของประตูหรือโต๊ะ
และแม้แต่ข้อมูลเหล่านี้ก็ขึ้นกับความยาวโฟกัส เพราะวัตถุเดียวกันที่ระยะเดียวกันจะดูใหญ่ขึ้นเมื่อใช้เลนส์ที่ยาวขึ้น
Eigen และคณะ~\cite{eigen14} เป็นงานแรก ๆ ที่ใช้โครงข่ายเชิงลึกกับปัญหานี้
และเพราะความกำกวมดังกล่าว จึงวัดความคลาดเคลื่อนบนลอการิทึมของความลึก
ในรูปแบบที่แยกความคลาดเคลื่อนของสเกลโดยรวมของฉากออกจากความคลาดเคลื่อนของรูปทรง
ฟังก์ชันสูญเสีย (loss) แบบไม่ขึ้นกับสเกลบนลอการิทึมนี้ในรูปแบบถ่วงน้ำหนัก เรียกว่า SiLog
ใช้ฝึกทั้ง ZoeDepth~\cite{zoedepth} และโมเดลที่ปรับจูนในงานนี้

โมเดลแบ่งได้เป็นสองกลุ่มตามวิธีจัดการกับสเกลที่หายไป กลุ่มแรกคือโมเดลความลึกสัมพัทธ์ (relative) ซึ่งไม่พยายามหาสเกล
MiDaS~\cite{midas} ฝึกโครงข่ายเดียวด้วยชุดข้อมูลหลายชุดที่วัดความลึกด้วยเซนเซอร์และหน่วยต่างกัน
ซึ่งทำได้ก็ต่อเมื่อฟังก์ชันสูญเสียไม่สนใจสเกลและค่าเลื่อน (offset) ของแต่ละภาพ
ผลลัพธ์จึงเป็นส่วนกลับของความลึก (disparity) ที่ถูกต้องเฉพาะเมื่อรู้สเกลและค่าเลื่อนของภาพนั้น
ต้องหาตัวเลขสองตัวนี้จากข้อมูลอื่นก่อนจึงจะได้ความลึกเป็นเมตร Depth Anything~\cite{da1}
และ Depth Anything~V2~\cite{da2} (\mbox{DA-v2}) ให้ผลลัพธ์แบบเดียวกัน แต่ฝึกด้วยภาพจำนวนมากกว่ามาก
กลุ่มที่สองคือโมเดลความลึกแบบเมตริก (metric) ซึ่งให้ความลึกเป็นเมตรโดยตรงโดยเรียนรู้สเกลจากข้อมูลที่มีค่าความลึกกำกับ
ZoeDepth~\cite{zoedepth} ต่อส่วนทำนายแบบเมตริกที่ปรับจูนด้วย NYU Depth~v2~\cite{nyuv2} และ KITTI~\cite{kitti}
เข้ากับโมเดลความลึกสัมพัทธ์ ส่วน \mbox{DA-v2} มีรุ่นเมตริกที่ปรับจูนจากรุ่นสัมพัทธ์ของตนเอง
โดยรุ่นในอาคาร (\mbox{DA-v2} Metric-Indoor) ปรับจูนด้วยภาพห้องที่เรนเดอร์จากชุดข้อมูล Hypersim~\cite{hypersim}
อีกแนวทางหนึ่งจัดการกับความยาวโฟกัสโดยตรง Metric3D~\cite{metric3d} แปลงทุกภาพให้เป็นภาพจากกล้องมาตรฐานที่มีความยาวโฟกัสคงที่
UniDepth~\cite{unidepth} ทำนายพารามิเตอร์ของกล้องไปพร้อมกับความลึก และ Depth Pro~\cite{depthpro}
ประมาณความยาวโฟกัสเองเมื่อไม่ได้รับค่านี้มา งานเหล่านี้รายงานผลเป็นตัวชี้วัดรายภาพบนชุดทดสอบมาตรฐาน
เช่น NYU Depth~v2 และ KITTI ตัวชี้วัดดังกล่าวเฉลี่ยจากภาพที่เป็นอิสระต่อกัน จึงไม่บอกว่าสเกลของโมเดลคงที่หรือไม่
ตลอดเฟรมของการสแกนครั้งเดียว ซึ่งเป็นสิ่งสำคัญเมื่อนำเฟรมทั้งหมดมารวมเป็นห้องเดียว

\subsection{โครงสร้างและการฝึกของ Depth Anything V2}
\mbox{DA-v2} ประกอบด้วยส่วนเข้ารหัสภาพ (encoder) และส่วนถอดรหัสความลึก (decoder)
ส่วนเข้ารหัสคือ DINOv2~\cite{dinov2} ซึ่งเป็น vision transformer (ViT)~\cite{vit}
โดยแบ่งภาพเป็นชิ้นสี่เหลี่ยมเล็ก ๆ (patch) แล้วให้แต่ละชิ้นดึงข้อมูลจากชิ้นอื่นทุกชิ้นผ่านกลไกที่เรียกว่า self-attention
DINOv2 ผ่านการฝึกล่วงหน้าโดยไม่ใช้ข้อมูลกำกับ บนชุดภาพขนาดใหญ่ที่คัดสรรแล้ว
เพื่อให้คุณลักษณะ (feature) ที่ได้ใช้กับงานได้หลากหลาย ส่วนถอดรหัสคือ Dense Prediction Transformer (DPT)~\cite{dpt}
ซึ่งนำคุณลักษณะจากหลายชั้นของส่วนเข้ารหัสมาประกอบกลับเป็นแผนที่คล้ายภาพที่หลายความละเอียด
แล้วรวมเป็นแผนที่ความลึกแบบหนาแน่นเพียงแผนที่เดียว โค้ดของโมเดลแบ่งส่วนถอดรหัสเป็น neck
ซึ่งทำหน้าที่ประกอบและรวมคุณลักษณะ กับ head ซึ่งให้ค่าความลึกออกมา ในรายงานนี้คำว่า DPT head
หมายถึงทั้งสองส่วน เช่นเดียวกับในบทนำ

\mbox{DA-v2} เองก็ฝึกด้วยครู (teacher) กล่าวคือ ฝึกโมเดลขนาดใหญ่ด้วยภาพสังเคราะห์ซึ่งมีความลึกที่แม่นยำก่อน
จากนั้นให้โมเดลนี้กำกับความลึกให้ภาพจริงที่ไม่มีความลึกราว 62 ล้านภาพ แล้วจึงให้โมเดลที่เผยแพร่เรียนรู้จากค่ากำกับที่โมเดลสร้างขึ้นนี้
ซึ่งเรียกว่า pseudo-label~\cite{da2} ผู้พัฒนา DINOv2 รายงานว่าทำนายความลึกได้ดีแม้ตรึงส่วนเข้ารหัสไว้
และฝึกเฉพาะส่วนถอดรหัส DPT ที่ต่อด้านบน~\cite{dinov2} การนำโมเดลที่ฝึกกับงานหรือโดเมนหนึ่งมาปรับใช้กับอีกงานหรือโดเมนหนึ่ง
เช่นนี้เรียกว่าการเรียนรู้แบบถ่ายโอน (transfer learning)~\cite{yosinski}
การฝึกเฉพาะส่วนถอดรหัสแบบที่งานนี้ทำใช้หน่วยความจำน้อยกว่ามาก
และยังคงคุณลักษณะทั่วไปที่ส่วนเข้ารหัสเรียนรู้มาจากภาพหลายล้านภาพไว้

\subsection{การฝึกด้วยความลึกจากเซนเซอร์}
โมเดลความลึกแบบมีผู้สอน (supervised) ใช้ความลึกที่วัดจากเซนเซอร์เป็นข้อมูลสอนมานานแล้ว โดยไม่ต้องให้คนกำกับ
ตัวอย่างเช่น NYU Depth~v2 ให้ความลึกในอาคารจากกล้อง Microsoft Kinect~\cite{nyuv2} KITTI ให้ความลึกบนถนนจากเครื่องสแกนเลเซอร์ที่ติดบนรถยนต์~\cite{kitti}
และ Eigen และคณะ~\cite{eigen14} ฝึกด้วยทั้งสองชุด สิ่งที่ต่างกันระหว่างงานเหล่านี้คือคุณภาพของเซนเซอร์
บน iPhone~Pro และ iPad~Pro นั้น ARKit รวมค่าที่วัดได้จาก LiDAR กับภาพสีเป็นแผนที่ความลึกขนาด
$256\times192$ พิกเซล พร้อมค่าความเชื่อมั่นของทุกพิกเซล~\cite{arkitscenes}
งานที่ทดสอบ LiDAR นี้เทียบกับวิธีสำรวจที่เป็นที่ยอมรับพบว่า แม่นยำราว 1~ซม. กับวัตถุที่ใหญ่กว่า 10~ซม.~\cite{luetzenburg}
และเพียงพอสำหรับการทำแผนที่สถาปัตยกรรมอย่างรวดเร็วที่มาตราส่วน 1:200~\cite{spreafico}
จึงมีเหตุผลที่จะใช้เป็นครู ARKitScenes~\cite{arkitscenes} ให้ความลึกนี้คู่กับความลึกจากเครื่องสแกนเลเซอร์ระดับงานสำรวจ
Faro~\cite{faro} ของห้องเดียวกัน โดยจัดให้ตรงกับเฟรมกล้องเดียวกัน และใช้ความลึกจาก Faro เป็นค่าอ้างอิง
ในโจทย์ที่เพิ่มความละเอียดของความลึกจาก LiDAR โดยอาศัยภาพสี
Prompt Depth Anything~\cite{promptda} ซึ่งเป็นงานที่ใช้ข้อมูลใกล้เคียงกับงานนี้ที่สุด ก็ป้อนความลึกจาก LiDAR ของ iPhone
เข้าโมเดล Depth Anything เป็นข้อมูลนำเข้าเพิ่มเติมเช่นกัน ในทั้งสองงาน LiDAR ต้องมีอยู่ทุกครั้งที่โมเดลทำงาน
ส่วนงานนี้ใช้ LiDAR เฉพาะตอนฝึก โมเดลที่นำไปใช้จึงต้องการเพียงภาพสี คำถามว่าครูที่มีสัญญาณรบกวนมากกว่านี้
ช่วยโมเดลขนาดใหญ่ที่ผ่านการฝึกล่วงหน้าได้เท่ากับเครื่องสแกนเลเซอร์หรือไม่
เมื่อตัดสินทั้งคู่ด้วยค่าอ้างอิงจากเลเซอร์ที่เป็นอิสระ คือเรื่องของวัตถุประสงค์ข้อ (1) และ (2)

\subsection{จากแผนที่ความลึกสู่แบบจำลองห้อง}
การรวมเฟรมต้องรู้ท่าทางกล้อง (camera pose) ของแต่ละเฟรม คือตำแหน่งและทิศทางของกล้อง
ARKit ติดตามท่าทางกล้องด้วย visual-inertial odometry (VIO) ซึ่งติดตามจุดเด่นในภาพจากเฟรมหนึ่งไปอีกเฟรมหนึ่ง
แล้วรวมกับเซนเซอร์วัดการเคลื่อนที่ของอุปกรณ์~\cite{vinsmono} เนื่องจากเซนเซอร์วัดการเคลื่อนที่วัดความเร่งเป็นหน่วยเมตริก
จุดที่ติดตามได้จึงมีตำแหน่งเป็นเมตร จุดเบาบาง (sparse points) เหล่านี้จึงให้สเกลและค่าเลื่อนที่โมเดลความลึกสัมพัทธ์ขาดไปได้ทีละเฟรม
แนวคิดการใช้ความลึกจากโครงข่ายร่วมกับเรขาคณิตแบบเบาบางของระบบติดตามมีมาตั้งแต่ CNN-SLAM~\cite{cnnslam}
และ MonoSDF~\cite{monosdf} ก็หาสเกลและค่าเลื่อนของทุกภาพก่อนนำความลึกจากภาพเดี่ยวไปช่วยสร้างพื้นผิว
การปรับสเกลด้วยจุดเบาบางของงานนี้ในหัวข้อ~\ref{sec:align} ก็อยู่ในแนวทางนี้

ขั้นต่อมาคือการสร้างห้อง ซึ่งทำโดยรวมความลึกเป็น truncated signed distance function (TSDF)~\cite{curless96}
โดยแบ่งพื้นที่เป็นลูกบาศก์เล็ก ๆ เรียกว่า voxel แต่ละ voxel เก็บระยะถึงพื้นผิวที่ใกล้ที่สุดที่สังเกตได้
เป็นบวกเมื่ออยู่หน้าพื้นผิวและเป็นลบเมื่ออยู่หลังพื้นผิว ตัดทิ้ง (truncate) เมื่อไกลเกินไม่กี่ voxel
แล้วเฉลี่ยค่านี้จากทุกเฟรมที่มองเห็น voxel นั้น จากนั้นจึงสกัดพื้นผิวตรงตำแหน่งที่ค่าเฉลี่ยเปลี่ยนเครื่องหมาย
ด้วยอัลกอริทึม marching cubes~\cite{marchingcubes} KinectFusion~\cite{kinectfusion} ทำให้การรวมแบบนี้ทำงานได้แบบเรียลไทม์
กับความลึกจาก Kinect และงานนี้ใช้การทำงานที่มีอยู่ใน Open3D~\cite{open3d}
การเฉลี่ยช่วยลดสัญญาณรบกวนที่สุ่มเปลี่ยนไปในแต่ละเฟรม แต่เฟรมที่มีสเกลไม่ตรงกันจะวางผนังเดียวกันไว้ที่ระยะต่างกัน
ค่าเฉลี่ยจึงทำให้ผนังเบลอหรือหนาขึ้น ความคงที่ของสเกลระหว่างเฟรม ซึ่งตัวชี้วัดรายภาพไม่ได้วัด จึงกำหนดคุณภาพของห้อง
นอกจากนี้ การรวมแบบ TSDF ไม่ขึ้นกับว่าความลึกเป็นเมตรนั้นมาจากแหล่งใด
จึงเป็นขั้นตอนคงที่ที่ยุติธรรมกับแหล่งความลึกทุกแหล่ง

ส่วนระบบสร้างแบบจำลองที่เรียนรู้ทั้งกระบวนการ (learned reconstruction) จะแทนที่บางขั้นตอนข้างต้น
ด้วยโครงข่ายประสาทเทียม Atlas~\cite{atlas} ทำนาย TSDF โดยตรง จากคุณลักษณะของภาพหลายภาพที่รู้ท่าทางกล้อง
NeuralRecon~\cite{neuralrecon}
ทำแบบเดียวกันทีละส่วนและแบบเรียลไทม์ และ SimpleRecon~\cite{simplerecon} ประมาณความลึกจากหลายเฟรมข้างเคียง
ก่อนรวมด้วยวิธีดั้งเดิม ทุกงานประเมินผลบน ScanNet~\cite{scannet} ซึ่งเป็นชุดการสแกนภาพสีและความลึกในอาคารขนาดใหญ่
แต่ละงานต้องใช้หลายเฟรมต่อการทำนายหนึ่งครั้ง และใช้โครงข่ายที่ฝึกมาสำหรับกระบวนการของตนเอง
จึงตอบคำถามต่างจากงานนี้ ซึ่งถามว่าโมเดลจากภาพเดี่ยวให้ห้องแบบใด เมื่อส่วนอื่นของกระบวนการไม่เปลี่ยน
งานนี้จึงไม่เปรียบเทียบกับงานกลุ่มนี้

\subsection{การวัดคุณภาพของแบบจำลองห้อง}
แบบจำลองห้องที่สร้างได้จะเทียบกับพื้นผิวอ้างอิงผ่านระยะระหว่างจุดที่สุ่มบนพื้นผิวทั้งสอง
Chamfer distance~\cite{chamfer} เฉลี่ยระยะจากแต่ละจุดถึงจุดที่ใกล้ที่สุดของอีกพื้นผิวทั้งสองทิศทาง
ส่วน Tanks and Temples~\cite{tanks} กำหนดระยะเกณฑ์ไว้ แล้วรายงานความแม่นยำ (precision)
คือสัดส่วนของพื้นผิวที่สร้างได้ที่อยู่ภายในระยะเกณฑ์จากพื้นผิวอ้างอิง ความครบถ้วน (recall)
คือสัดส่วนของพื้นผิวอ้างอิงที่อยู่ภายในระยะเกณฑ์จากพื้นผิวที่สร้างได้ และค่าเฉลี่ยฮาร์มอนิกของทั้งสอง คือ F-score
ระยะเกณฑ์ทำให้คะแนนผูกกับการใช้งาน ซึ่งเป็นวิธีที่บทนำใช้เชื่อมระยะ 2, 5 และ 10~ซม.
กับงานวัดเพื่อปรับปรุงห้อง งานวางเฟอร์นิเจอร์ และงานดูภาพรวมของห้อง สำหรับการวัดรายเฟรม
งานนี้ใช้ AbsRel และ \wone{} จากงานด้านการประมาณความลึกจากภาพเดี่ยว~\cite{eigen14}

\subsection{ตำแหน่งของงานวิจัยนี้}
โดยสรุป งานก่อนหน้ารายงานความแม่นยำรายภาพของโมเดลความลึกบนชุดทดสอบมาตรฐาน
ปรับโมเดลเมตริกด้วยข้อมูลกำกับจากภาพสังเคราะห์หรือจากเซนเซอร์ของชุดข้อมูลสาธารณะ
ป้อน LiDAR ของอุปกรณ์พกพาเข้าโมเดลขณะใช้งาน หรือทดสอบตัว LiDAR ของอุปกรณ์พกพาเองเทียบกับวิธีสำรวจ
เราไม่พบงานที่รวมประเด็นเหล่านี้ไว้แบบงานนี้ กล่าวคือ ใช้ LiDAR ของอุปกรณ์ผู้บริโภคเป็นครูเฉพาะตอนฝึก
เปรียบเทียบกับการใช้เครื่องสแกนเลเซอร์ระดับงานสำรวจเป็นครูบนห้องชุดเดียวกัน
และตัดสินทั้งคู่ด้วยค่าอ้างอิงจากเลเซอร์ที่เป็นอิสระ ทั้งเป็นเซนติเมตรบนแบบจำลองห้องและแบบรายเฟรม


\bibliographystyle{IEEEtran}
\bibliography{refs}

\end{document}
